\section{El espacio de riesgo-recompensa de los proyectos de IA}

El ponente toma prestados dos conceptos clásicos de la gestión de inversiones —\textbf{Risk} (riesgo) y \textbf{Reward} (recompensa)— para construir un marco bidimensional que permite analizar los proyectos de IA.

\subsection{Matriz de riesgo-recompensa}

En este marco, el eje horizontal representa Reward (recompensa/utilidad) y el eje vertical representa Risk (riesgo/daño potencial). El ponente interactúa con los estudiantes para ubicar distintas aplicaciones de IA en este espacio:

\begin{table}[H]
\centering
\caption{Ubicación de las aplicaciones de IA en el espacio de riesgo-recompensa}
\begin{tabular}{@{}lccl@{}}
\toprule
\textbf{Aplicación de IA} & \textbf{Risk} & \textbf{Reward} & \textbf{Descripción} \\
\midrule
Code Completion (una sola línea) & Bajo & Bajo & Completa la firma de una función; un error tiene poco impacto \\
Text Revision (revisión de texto) & Bajo--Medio & Medio & Devuelve una versión mejorada del texto; el riesgo es un poco más alto \\
Text Summary (resumen de texto) & Medio & Medio--Alto & Puede omitir o malinterpretar información clave \\
LLM (modelo de lenguaje grande) & Medio--Alto & Alto & Tiene capacidades asombrosas, pero conlleva varios riesgos \\
AI Judge (juez automático) & Muy alto & Incierto & Involucra la libertad y la vida de las personas \\
AGI (inteligencia artificial general) & Extremo & Extremo & Una recompensa alta viene acompañada de un riesgo incontrolable \\
\bottomrule
\end{tabular}
\end{table}

\begin{warningbox}{La doble naturaleza extrema de la AGI}
El ponente ubica la AGI en la esquina superior derecha del espacio de riesgo-recompensa: tiene tanto una recompensa potencial extrema (una inteligencia general podría resolver casi todos los problemas que enfrenta la humanidad) como un riesgo extremo (comportamiento impredecible, posible pérdida de control). Esta combinación de «alto riesgo y alta recompensa» convierte a la AGI en el objetivo de IA más controvertido. Como en la escena de la película \textit{Terminator}, si una superinteligencia llega a perder el control, las consecuencias serán irreversibles.
\end{warningbox}

\subsection{Umbral de riesgo: ¿dónde trazar la línea?}

El ponente plantea una pregunta clave: \textbf{¿en qué punto del eje de riesgo deberíamos trazar una línea roja «infranqueable»?}

\begin{itemize}
    \item \textbf{Estrategia conservadora}: la línea roja se traza muy baja y solo se permiten aplicaciones de bajo riesgo (como la completación de código o la revisión de texto)
    \item \textbf{Estrategia agresiva}: la línea roja se traza muy alta, se tolera la mayor parte del riesgo y solo se prohíben los escenarios extremos
    \item \textbf{Dilema real}: las distintas partes interesadas tienen juicios muy distintos sobre el umbral de riesgo
\end{itemize}

Un estudiante mencionó una tercera dimensión importante: \textbf{la capacidad económica}. Las grandes empresas cuentan con equipos legales suficientes para afrontar las consecuencias del riesgo, mientras que las empresas emergentes pueden no tener la capacidad de asumir el costo de las disputas legales. Esta asimetría hace que el problema de «trazar la línea» sea aún más complejo.

\begin{knowledgebox}{La ley solo se preocupa por el riesgo, no por la recompensa}
El ponente señala una distinción importante: el sistema legal \textbf{no toma en cuenta la dimensión de la recompensa}. No se puede argumentar ante un tribunal que «mi IA es muy útil, así que se le debería permitir infringir las reglas»: la ley solo se preocupa por si el sistema causó daño o violó alguna norma. Esto significa que, dentro del marco legal, el análisis de riesgo-recompensa se reduce a una evaluación puramente de riesgo.
\end{knowledgebox}

\subsection{Resumen de la sección}

El marco de riesgo-recompensa ofrece una herramienta intuitiva para pensar los proyectos de IA. Las distintas aplicaciones de IA ocupan posiciones diferentes en este espacio, y la pregunta de «dónde trazar la línea» es, en esencia, una decisión multidimensional que involucra aspectos técnicos, éticos, legales y económicos. El sistema legal tiende a centrarse únicamente en la dimensión del riesgo, mientras que los desarrolladores necesitan encontrar un punto de equilibrio entre el atractivo de la recompensa y las restricciones del riesgo.

